\documentclass[aspectratio=169, 11pt]{beamer}
\usepackage{fontspec}
\setmainfont{TeX Gyre Termes}
\setsansfont{TeX Gyre Heros}
\setmonofont{Latin Modern Mono}
\usepackage[spanish,es-nodecimaldot]{babel}
\usepackage{amsmath, amssymb, amsfonts, booktabs, graphicx}
\usepackage{tikz}
\usetikzlibrary{shapes.geometric, arrows.meta, positioning, calc, fit, shadows}
\usetheme{Madrid}
\usecolortheme{beaver}
\definecolor{unalRojo}{RGB}{148, 36, 36}
\definecolor{verdeOK}{RGB}{34, 139, 34}
\definecolor{naranjaAlerta}{RGB}{230, 126, 34}
\definecolor{rojoAlerta}{RGB}{192, 57, 43}
\setbeamercolor{structure}{fg=unalRojo}
\setbeamercolor{block title}{bg=unalRojo!15, fg=unalRojo}
\setbeamercolor{block body}{bg=unalRojo!5}
\setbeamercolor{alerted text}{fg=rojoAlerta}
\newcommand{\BSF}{\textit{Hermetia illucens}}
\newcommand{\MAPE}{\text{MAPE}}
\newcommand{\DEB}{\text{DEB}}
\newcommand{\NDIR}{\text{NDIR}}
\newcommand{\MRV}{\text{MRV}}
\setbeamertemplate{navigation symbols}{}
\setbeamertemplate{footline}[frame number]

\begin{document}

% ---------------------------------------------------------------------------
% Slide 1: Conservaci\'{o}n de masa
% ---------------------------------------------------------------------------
\begin{frame}{Conservaci\'{o}n de masa en el metabolismo larval}
    \begin{center}
        \begin{tikzpicture}[scale=0.85, transform shape,
            box/.style={rectangle, draw=unalRojo, fill=unalRojo!10, rounded corners, minimum width=2.4cm, minimum height=1cm, align=center, font=\small}]
            \node[box] (ing) at (0,0) {Ingesta};
            \node[box] (asi) at (3.5,0) {Asimilaci\'{o}n neta};
            \node[box] (bio) at (7,1.2) {Biomasa $B(t)$};
            \node[box] (res) at (7,-1.2) {Reservas};
            \node[box] (rsp) at (10.5,1.2) {Respiraci\'{o}n CO$_2$};
            \node[box] (exc) at (10.5,-1.2) {Excreci\'{o}n};

            \draw[-{Stealth}, thick, unalRojo] (ing) -- (asi);
            \draw[-{Stealth}, thick, unalRojo] (asi) -- (bio);
            \draw[-{Stealth}, thick, unalRojo] (asi) -- (res);
            \draw[-{Stealth}, thick, unalRojo] (bio) -- (rsp);
            \draw[-{Stealth}, thick, unalRojo] (res) -- (rsp);
            \draw[-{Stealth}, thick, unalRojo] (asi) -- (exc);
        \end{tikzpicture}
    \end{center}

    \vfill
    \begin{block}{Ecuaci\'{o}n diferencial ordinaria (Eriksen 2022)}
        \centering
        $\displaystyle \frac{dB}{dt} = \underbrace{\text{Asimilaci\'{o}n}}_{\text{ingesta } \times \text{ eficiencia}} - \underbrace{\text{Mantenimiento}}_{m \cdot B(t)} - \underbrace{\text{S\'{i}ntesis}}_{Y_{\text{CO}_2}^{-1} \cdot \frac{dB}{dt}}$
    \end{block}
\end{frame}

% ---------------------------------------------------------------------------
% Slide 2: Cin\'{e}tica acoplada al ciclo
% ---------------------------------------------------------------------------
\begin{frame}{Cin\'{e}tica acoplada al ciclo de bioconversi\'{o}n}
    \begin{columns}[T]
        \begin{column}{0.5\textwidth}
            \begin{itemize}
                \item Fase exponencial (d\'{i}as 0--10): $B(t)$ crece y la tasa de CO$_2$ se acelera.
                \item Punto de inflexi\'{o}n: la ingesta se satura por autocompetencia.
                \item Fase decreciente: la biomasa anticipa la prepupaci\'{o}n.
            \end{itemize}
        \end{column}
        \begin{column}{0.45\textwidth}
            \begin{tikzpicture}[scale=0.75, transform shape]
                \draw[-{Stealth}] (0,0) -- (6,0) node[right, font=\small] {Tiempo};
                \draw[-{Stealth}] (0,0) -- (0,4.2) node[above, font=\small] {Tasa CO$_2$};
                \draw[very thick, unalRojo] plot[smooth, tension=0.8] coordinates {(0.3,0.4) (1.2,1.2) (2.5,3) (3.8,3.3) (4.6,2.2) (5.5,0.8)};
                \node[font=\footnotesize, naranjaAlerta] at (1.5,-0.5) {Exponencial};
                \node[font=\footnotesize, naranjaAlerta] at (3.8,-0.5) {M\'{a}ximo};
                \node[font=\footnotesize, naranjaAlerta] at (5.2,-0.5) {Decreciente};
            \end{tikzpicture}
        \end{column}
    \end{columns}
\end{frame}

\end{document}
